\documentclass[11pt,a4paper]{article}

% ======================================================================
%  Technical Appendix: Heat Kernel, Fermions, and the Sign of Induced
%  Gravity. Version 2 (July 2026). v1: December 2025 (ai.viXra:2512.0091v1).
%  This is a self-contained heat-kernel / induced-gravity support piece,
%  topically adjacent to (not part of) the coherence-maintenance series.
%  Main changes (Appendix): the physics is UNCHANGED and correct -- v1 had
%  no error; this v2 (i) updates the companion references to their v2
%  revisions, (ii) adds a symbolic verification of the full Seeley-DeWitt
%  a1 coefficient table and the induced-Newton sign column (all seven
%  species checked, distributed as verification/0091/verify_0091.py), and
%  (iii) tightens the scope note so the link to the "resource boundary /
%  cut" program is stated as heuristic, not load-bearing.
% ======================================================================

\usepackage[utf8]{inputenc}
\usepackage[T1]{fontenc}
\usepackage{amsmath,amssymb}
\usepackage[margin=1.05in]{geometry}
\usepackage{booktabs}
\usepackage[colorlinks=true,linkcolor=blue,citecolor=blue,urlcolor=blue]{hyperref}

\newcommand{\Aeff}{A_1^{(\mathrm{eff})}}

\title{Technical Appendix: Heat Kernel, Fermions, and the Sign of Induced
Gravity\\[2mm]
\large (Sign conventions fixed; Laplacian/Lichnerowicz hinge made explicit)}
\author{Lluis Eriksson\\ \small Independent Researcher\\
\small \texttt{lluiseriksson@gmail.com}}
\date{July 4, 2026 \quad (v2; v1: December 2025) \quad ai.viXra:2512.0091}

\begin{document}

\maketitle

\begin{center}
\fbox{\parbox{0.94\textwidth}{\small
\textbf{Scope note.} This appendix is a self-contained heat-kernel
computation of the sign of the induced Einstein--Hilbert term from integrating
out matter fields. It is \emph{topically adjacent} to the coherence-
maintenance / ``resource boundary'' program (the induced-gravity sign is a
curvature-sector bookkeeping that the program refers to heuristically), but it
is \emph{not} load-bearing for any result in that series and depends on none
of it; the cross-references \cite{GML,MC,QCB,Davies,RIP,Closure,Stress,HCut,Work,Clustering}
are for context only. It does not attempt a UV completion nor a physical
determination of $G_{\mathrm{phys}}$. \textbf{v2}: the physics is unchanged
(v1 contained no error); this revision updates companion references to their
v2 revisions and adds a symbolic verification of the full coefficient table
(Appendix~\ref{app:verif}).}}
\end{center}

\begin{abstract}
We derive the local contribution proportional to the scalar curvature $R$ in
the Euclidean one-loop effective action obtained by integrating out matter
fields on a curved background. Using a Schwinger cutoff $\varepsilon =
\Lambda^{-2}$ and the Seeley--DeWitt coefficient $a_1$, we extract the
quadratically divergent term multiplying $\int d^4x\sqrt{g}\,R$. We fix a
single Euclidean convention for the Einstein--Hilbert action, state an explicit
Laplacian convention, and write the Lichnerowicz/Weitzenb\"ock identity in a
sign-robust form so that the fermionic contribution is unambiguous. We provide
a unified bookkeeping coefficient $\Aeff$ such that
\[
W_R^{\mathrm{total}} = -\frac{\Aeff}{32\pi^2}\Lambda^2
\int d^4x\sqrt{g}\,R,
\]
and hence an induced Newton constant $G_{\mathrm{ind}}$ via comparison with the
Euclidean Einstein--Hilbert action. We include the minimal gauge${+}$ghost
package in background Feynman gauge, a species table, and a reproducible
symbolic check of every entry.
\end{abstract}

\section{Conventions and notation}

\subsection{Lorentzian geometry and curvature sign}
Lorentzian signature $(-,+,+,+)$ and the Wald/MTW convention
\[
R^a{}_{bcd} = \partial_c\Gamma^a_{bd} - \partial_d\Gamma^a_{bc}
+ \Gamma^a_{ce}\Gamma^e_{bd} - \Gamma^a_{de}\Gamma^e_{bc},
\quad R_{bd} = R^a{}_{bad}, \quad R = g^{bd}R_{bd}.
\]

\subsection{Euclidean formulation and Einstein--Hilbert action}
Heat-kernel calculations are done in $D=4$ Euclidean signature. We fix
\[
S_E^{\mathrm{EH}} = -\frac{1}{16\pi G}\int d^4x\sqrt{g}\,R,
\qquad
S_L = \frac{1}{16\pi G}\int d^4x\sqrt{-g}\,R,
\]
with $S_E = -iS_L$ when topological/eta phases are irrelevant for the local
divergent coefficient extracted here.

\subsection{Laplacian, Laplace-type operators, Schwinger cutoff}
$\nabla^2 \equiv g^{\mu\nu}\nabla_\mu\nabla_\nu$, so $-\nabla^2$ is positive; a
Laplace-type operator is $P = -\nabla^2 + X(x)$ with $X$ a local endomorphism.
Schwinger proper-time with $\varepsilon = \Lambda^{-2}$:
\[
\ln\det P = -\int_\varepsilon^\infty \frac{ds}{s}\,\mathrm{Tr}\,e^{-sP}.
\]

\subsection{Fermionic hinge: Lichnerowicz/Weitzenb\"ock}
$D_E \equiv \gamma^\mu\nabla_\mu$ with Euclidean $(\gamma^\mu)^\dagger =
\gamma^\mu$ and $(D_E)^\dagger = -D_E$;
$[\nabla_\mu,\nabla_\nu]\psi = \tfrac14 R_{\mu\nu\rho\sigma}\gamma^{\rho\sigma}
\psi$, $\gamma^{\rho\sigma} = \tfrac12[\gamma^\rho,\gamma^\sigma]$. In
sign-robust form,
\[
-D_E^2 = -\nabla^2 + \tfrac14 R.
\]
Here $\nabla^2$ is the connection Laplacian on spinors (not the de Rham
Laplacian on forms).

\section{Heat kernel expansion and $a_1$}
With $K(s;x,y) = \langle x|e^{-sP}|y\rangle$,
\[
\mathrm{Tr}\,e^{-sP} \sim \frac{1}{(4\pi s)^2}\int d^4x\sqrt{g}
\big(a_0 + a_1 s + a_2 s^2 + \cdots\big),
\quad
a_0 = \mathrm{tr}\,\mathbf{1},\;\;
a_1 = \mathrm{tr}\big(\tfrac16 R - X\big).
\]
For a real scalar with $X = m^2 + \xi R$: $a_{1,R} = \tilde a_1 R$,
$\tilde a_1 = \tfrac16 - \xi$; the $-m^2$ piece renormalizes the cosmological
term, not $\int\sqrt{g}\,R$.

\section{Scalars: extraction of $\int\sqrt{g}\,R$}
$W_{\mathrm{scalar}} = \tfrac12\ln\det P = -\tfrac12\int_\varepsilon^\infty
\tfrac{ds}{s}\mathrm{Tr}\,e^{-sP}$; using $\int_\varepsilon^\infty ds\,s^{-2}
= \Lambda^2$,
\[
W_R^{\mathrm{scalar}} = -\frac{\tilde a_1}{32\pi^2}\Lambda^2
\int d^4x\sqrt{g}\,R.
\]

\section{Dirac fermions: reduction to Laplace type}
$S_F = \int d^4x\sqrt{g}\,\bar\psi(D_E + m)\psi$, $W_{\mathrm{ferm}} =
-\ln\det(D_E + m)$. With $\ln\det(D_E+m) = \tfrac12\ln\det[(D_E+m)^\dagger
(D_E+m)] + i\Theta$ and $P_D := (D_E+m)^\dagger(D_E+m) = -D_E^2 + m^2 =
-\nabla^2 + m^2 + \tfrac14 R$, so $X_D = m^2 + \tfrac14 R$ and
\[
a_1^{(P_D)} = \tfrac16 R - X_D = -\tfrac{1}{12}R - m^2,
\qquad
\mathrm{tr}\,a_{1,R}^{(P_D)} = 4\cdot\big(-\tfrac{1}{12}\big)R = -\tfrac13 R.
\]
Because $W_{\mathrm{ferm}} = -\tfrac12\ln\det P_D = +\tfrac12
\int_\varepsilon^\infty\tfrac{ds}{s}\mathrm{Tr}\,e^{-sP_D}$ (the fermionic
sign flip),
\[
W_R^{\mathrm{ferm}} = -\frac{1}{96\pi^2}\Lambda^2\int d^4x\sqrt{g}\,R.
\]
The phase $i\Theta$ and zero modes matter in chiral/topological settings but do
not modify this local coefficient for a massive Dirac field in a smooth
background.

\section{Gauge fields and ghosts (spin-1), Feynman gauge}
Per generator, $\alpha=1$: $(P_1)_\mu{}^\nu = -\delta_\mu^\nu\nabla^2 +
R_\mu{}^\nu$, ghost $P_{\mathrm{gh}} = -\nabla^2$,
$W_{\mathrm{gauge}} = \tfrac12\ln\det P_1 - \ln\det P_{\mathrm{gh}}$. With
$E_\mu{}^\nu = R_\mu{}^\nu$,
\[
\mathrm{tr}\,a_{1,R}^{(1\text{-form})} = \tfrac16\,\mathrm{tr}\,\mathbf 1
- \mathrm{tr}\,E = \tfrac46 - 1 = -\tfrac13 R;
\]
the ghost (complex Grassmann scalar) enters with the opposite sign of a
complex boson ($\tfrac13 R$), giving $\mathrm{tr}\,a_{1,R}^{(\mathrm{gh})} =
-\tfrac13 R$, hence
\[
\Aeff(\text{vector}+\text{ghosts}) = -\tfrac23 \quad(\text{per generator}).
\]

\section{Unified convention and induced Newton constant}
Define $\Aeff$ by $W_R^{\mathrm{total}} = -\dfrac{\Aeff}{32\pi^2}\Lambda^2
\int d^4x\sqrt{g}\,R$. Comparing with $S_E^{\mathrm{EH}}$ and identifying
$S_{E,\mathrm{induced}}^{\mathrm{EH}} = W_R^{\mathrm{total}}$,
\[
\frac{1}{16\pi G_{\mathrm{ind}}} = \frac{\Aeff}{32\pi^2}\Lambda^2
\quad\Longrightarrow\quad
G_{\mathrm{ind}} = \frac{2\pi}{\Aeff\,\Lambda^2},
\qquad
G_{\mathrm{ind}} > 0 \iff \Aeff > 0.
\]

\section{Species table (unified convention)}
\begin{center}
\begin{tabular}{lcl}
\toprule
Species & $\Aeff$ & Comment\\
\midrule
Real scalar (general $\xi$) & $\tfrac16 - \xi$ & $X = m^2 + \xi R$\\
Real scalar (minimal) & $+\tfrac16$ & $\xi = 0$\\
Complex scalar & $2(\tfrac16 - \xi)$ & two real d.o.f.\\
Dirac fermion (4 comp.) & $+\tfrac13$ & derived above\\
Weyl fermion (2 comp.) & $+\tfrac16$ & half of Dirac\\
Majorana fermion & $+\tfrac16$ & half of Dirac\\
Gauge vector ${+}$ ghosts & $-\tfrac23$ & per generator, Feynman gauge\\
\bottomrule
\end{tabular}
\end{center}
Every entry is verified symbolically in Appendix~\ref{app:verif}.

\section{Sign examples}
\begin{itemize}
\item Minimal real scalar: $\Aeff = \tfrac16 > 0 \Rightarrow G_{\mathrm{ind}} >
0$.
\item Non-minimally coupled real scalar with $\xi = \tfrac14$: $\Aeff =
\tfrac16 - \tfrac14 = -\tfrac{1}{12} < 0 \Rightarrow G_{\mathrm{ind}} < 0$.
\item Conformally coupled real scalar in $D=4$ ($\xi_{\mathrm{conf}} =
\tfrac{D-2}{4(D-1)} = \tfrac16$): $\Aeff = 0$; it does not contribute to the
induced Einstein--Hilbert term at this order. This does not remove the usual
conformal-anomaly / $a_2$ curvature-squared terms; it only says that the
quadratically divergent Einstein--Hilbert $R$-term vanishes for the
conformally coupled scalar.
\item Gauge vector (per generator): $\Aeff = -\tfrac23 < 0 \Rightarrow
G_{\mathrm{ind}} < 0$ for that sector alone.
\end{itemize}
The net sign for a given model is $\sum_{\mathrm{species}} \Aeff$ weighted by
multiplicities; realistic spectra require summing all sectors (this appendix
provides the per-species bookkeeping, not a spectrum-specific determination).

\appendix

\section{Symbolic verification}\label{app:verif}

Distributed as \texttt{verification/0091/verify\_0091.py} (SymPy). It
reconstructs each $\Aeff$ from $a_1 = \mathrm{tr}(\tfrac16 R - X)$ with the
Lichnerowicz shift $X_D = m^2 + \tfrac14 R$ for spinors, the fermionic
overall sign flip ($W_{\mathrm{ferm}} = -\tfrac12\ln\det P_D$), and the ghost
entering as minus a complex boson, and checks all seven table entries:
$\tfrac16-\xi$, $\tfrac16$, $\tfrac13$, $\tfrac13$, $\tfrac16$, $\tfrac16$,
$-\tfrac23$ --- all match. The induced-Newton sign column ($G_{\mathrm{ind}}
> 0 \iff \Aeff > 0$) is checked on the minimal scalar ($+$), the non-minimally coupled scalar with
$\xi=\tfrac14$ ($-$), the conformally coupled scalar in $D=4$ ($\xi=\tfrac16$,
vanishing contribution), and the gauge sector ($-$). Result: \textsc{all Seeley--DeWitt coefficients
match}.

\section{Changes relative to version 1}

\begin{enumerate}
\item \textbf{Physics unchanged (no error in v1).} The coefficient table and
sign bookkeeping are correct and verified (Appendix~\ref{app:verif}); this is
the sole paper in the program requiring no mathematical rectification.
\item \textbf{Verification added:} symbolic check of the full $a_1$ table and
the induced-Newton sign, distributed with the paper.
\item \textbf{References updated} to the v2 revisions of the companion series;
the ``resource boundary / cut'' link is restated as heuristic and non-load-
bearing (scope note).
\end{enumerate}

\begin{thebibliography}{99}\small
\bibitem{GML} L.~Eriksson, \emph{Geometry, Membranes, and Life as a Resource
Boundary}, ai.viXra:2512.0085, v2 (2026; v1: 2025).
\bibitem{MC} L.~Eriksson, \emph{The Maintenance Constraint}, ai.viXra:2512.0084,
v2 (2026; v1: 2025).
\bibitem{QCB} L.~Eriksson, \emph{Operational Coherence Maintenance and the
Quantum--Classical Boundary}, ai.viXra:2512.0081, v2 (2026; v1: 2025).
\bibitem{Davies} L.~Eriksson, \emph{Finite-Dimensional Davies Interface Lemmas
and TFIM Witness Tests}, ai.viXra:2512.0073, v2 (2026; v1: 2025).
\bibitem{RIP} L.~Eriksson, \emph{The Rate Inheritance Principle},
ai.viXra:2512.0072, v2 (2026; v1: 2025).
\bibitem{Closure} L.~Eriksson, \emph{Operational Coherence Maintenance: Proven
Results, Conditional Interfaces, and Open Dynamical Gaps}, ai.viXra:2512.0071,
v2 (2026; v1: 2025).
\bibitem{Stress} L.~Eriksson, \emph{Stress Testing the Rate Inheritance
Principle}, ai.viXra:2512.0070, v2 (2026; v1: 2025).
\bibitem{HCut} L.~Eriksson, \emph{The Heisenberg Cut as a Resource Boundary},
ai.viXra:2512.0064, v2 (2026; v1: 2025).
\bibitem{Work} L.~Eriksson, \emph{The Conditional Maintenance Work Theorem},
ai.viXra:2512.0061, v2 (2026; v1: 2025).
\bibitem{Clustering} L.~Eriksson, \emph{Clustering, Recovery, and Locality in
AQFT}, ai.viXra:2512.0060, v2 (2026; v1: 2025).
\bibitem{Gilkey} P.~B.~Gilkey, \emph{Invariance Theory, the Heat Equation, and
the Atiyah--Singer Index Theorem}, CRC Press (1995).
\bibitem{Vassilevich} D.~V.~Vassilevich, \emph{Heat kernel expansion: user's
manual}, Phys. Rept. \textbf{388} (2003) 279, arXiv:hep-th/0306138.
\bibitem{Avramidi} I.~G.~Avramidi, \emph{Heat Kernel and Quantum Gravity},
Lecture Notes in Physics (2000).
\bibitem{Sakharov} A.~D.~Sakharov, \emph{Vacuum quantum fluctuations in curved
space and the theory of gravitation}, Dokl. Akad. Nauk SSSR \textbf{177}
(1967) 70.
\bibitem{BarvinskyVilkovisky} A.~O.~Barvinsky and G.~A.~Vilkovisky, \emph{The
generalized Schwinger--DeWitt technique}, Phys. Rept. \textbf{119} (1985) 1.
\end{thebibliography}

\end{document}
